\section{Related Work}
\label{sec:research}
\paragraph{ViZDoom as a research platform.}
\citet{kempka2016vizdoom} introduced \vzd{} as a platform for visual reinforcement learning, enabling controlled first-person tasks and access to game-derived observations. \citet{wydmuch2019competitions} subsequently analyzed agents and evaluation in multiplayer competitions. We use this latter work as the historical reference: its preprint appeared in 2018, followed by the journal publication in 2019. Depth observations, instance labels, map access, and native multiplayer were already part of the platform's earlier scope. The present contribution extends that scope through additional modalities, semantic information, standardized task interfaces, and language-model interaction.

\paragraph{Prediction, control, and memory.}
Several influential studies used \vzd{} to connect perception to decisions over time. In \emph{Learning to Act by Predicting the Future}, \citet{dosovitskiy2017future} selected actions using predicted future measurements, enabling control conditioned on different goals. \citet{lample2017fps} combined recurrent decision-making with auxiliary game-feature prediction for first-person combat. These works established Doom as a setting in which temporally informed representations can support control beyond individual image recognition.

\paragraph{Exploration and world models.}
\citet{pathak2017curiosity} evaluated curiosity-driven exploration in \vzd{}, using prediction error to encourage behavior when external rewards provide little guidance. \citet{savinov2019episodic} investigated episodic curiosity based on reachability and memory, including Doom navigation experiments. \citet{ha2018worldmodels} used the Take Cover scenario to study a controller trained through learned latent dynamics. Together, these studies motivate evaluating exploration, remembered experience, and predictive representations under longer and more varied task dependencies.

\paragraph{Generalization, continual learning, and safety.}
MazeExplorer introduced procedurally generated \vzd{} mazes with separate training and test environments~\citep{harries2019mazeexplorer}. LevDoom studied generalization across levels of environmental difficulty~\citep{tomilin2022levdoom}, while COOM used varied Doom domains and objectives to examine continual reinforcement learning~\citep{tomilin2023coom}. HASARD extended this line of work to vision-based safe reinforcement learning with explicit task rewards and costs~\citep{tomilin2025hasard}. These publications illustrate the value of adapting a shared engine to distinct research questions, and the need to evaluate more than performance on a single familiar scenario.

\paragraph{Learning systems and multi-agent interfaces.}
Sample Factory demonstrated efficient learning and multiplayer self-play in Doom~\citep{petrenko2020samplefactory}, showing how learning infrastructure and agent populations contribute to effective visual control. PettingZoo provides shared interfaces for multi-agent environments~\citep{terry2021pettingzoo}. Integrating \vzd{} with this ecosystem makes its multiplayer tasks more accessible to established multi-agent methods while retaining the game's partially observable interaction setting.

\paragraph{Contemporary agent evaluation.}
Broader research provides context for the proposed benchmark. DreamerV3 demonstrates world-model learning across diverse control domains~\citep{hafner2025dreamer}; BALROG studies language and vision-language models through interactive games~\citep{paglieri2025balrog}; and SIMA~2 investigates generalist agents in virtual worlds~\citep{sima2025}. We cite these as methodological context, rather than as additional examples of \vzd{} use. Procgen emphasizes evaluation over environment distributions~\citep{cobbe2020procgen}, and \citet{agarwal2021precipice} motivate reporting uncertainty and robust aggregate performance instead of relying on isolated point estimates.

MCP supplies a shared protocol through which AI applications access tools and external context~\citep{mcprotocol}. Harbor provides infrastructure for evaluating models and agent systems in containerized tasks~\citep{harborframework}. Our contribution is the \vzd{} interaction server and the associated task and evaluation design. By combining controlled scenarios, full-level objectives, and a common tool interface, the benchmark can examine where failures arise from individual skills, their composition, or the harness that coordinates them.
